\section*{Заключение}
\addcontentsline{toc}{section}{Заключение}

В ходе выполнения лабораторной работы была изучена полиномиальная аппроксимация
методом наименьших квадратов. По таблице из $n = 20$ узлов были построены два
аппроксимирующих полинома.

Для линейной аппроксимации $P_1(x) = a_0 + a_1 x$ нормальная система МНК
была составлена и решена, в результате чего получены коэффициенты
$a_0 = 3{,}4448$ и $a_1 = -1{,}0327$:
\[
P_1(x) = 3{,}4448 - 1{,}0327\,x.
\]
Наилучшее среднеквадратическое отклонение составило $\delta_2(P_1) = 0{,}1464$.

Для квадратичной аппроксимации $P_2(x) = a_0 + a_1 x + a_2 x^2$ получены
коэффициенты $a_0 = 3{,}1147$, $a_1 = -0{,}1323$, $a_2 = -0{,}4287$:
\[
P_2(x) = 3{,}1147 - 0{,}1323\,x - 0{,}4287\,x^2.
\]
Наилучшее среднеквадратическое отклонение составило $\delta_2(P_2) = 0{,}0728$.

Квадратичная модель вдвое точнее линейной: $\delta_2(P_2) < \delta_2(P_1)$.
Это закономерно, поскольку полином более высокой степени обладает большей
гибкостью при подгонке к данным. Все вычисления выполнены программно на языке
Python с использованием библиотеки NumPy.